\section{Data generation}
\label{sec:coral-island-dataset-generation}

The creation of the dataset is done through the use of the procedural algorithm for which we alter the input parameters.

For each generation, the top-view and profile-view sketches use the user-given input. Each outline of the top-view sketch $\radius$ is modified by adding continuous noise $\noise$, giving the final outlines:
\begin{align*}
    \radius^*(\angl) = \radius(\angl) + \noise_{\text{outlines}}(\angl).
\end{align*}

On the other hand, we define a new initial profile-view sketch by defining $\heightProfile^*(\distRegions_\p)$, the initial height function on which fBm noise $\noise_{\text{height}}$ is applied to obtain
\begin{align*}
    \heightProfile^*(\distRegions_\p) = \heightProfile(\distRegions_\p) \cdot \noise_{\text{height}}(\distRegions_\p).
\end{align*}

An identical process is done for the resistance function:
\begin{align*}
    \resistance^*(\distRegions_\p) = \resistance(\distRegions_\p) \cdot \noise_{\text{resistance}}(\distRegions_\p).
\end{align*}

Next, we generate a random deformation field. We do not target realism in stroke generation. Instead we generate a random number $n$ of strokes with random paths. Spread and intensity of each stroke is also randomised, providing complexity and variety in the results. Our random wind stroke paths are smooth, simplex noise-driven trajectories biased toward the island's center, ensuring that distortions remain mainly around the main landmass.

Finally, to further enrich the dataset with erosion-like structures, we add a set of radial strokes of small width and low strength, extending between the island centre and the exterior of the canvas, oriented alternately landward and seaward. This produces simple radial drainage patterns illustrated in \cref{fig:coral-island-radial-pattern-example}, available by the radial definition of our islands contours.

\begin{figure}[tb]
    \autofitgraphics[]{circular-without-radial-erosion.png, circular-without-radial-erosion-render-wireframe.png, circular-with-radial-erosion.png, circular-with-radial-erosion-render-wireframe.png}
    \autofitcaptions[]{Height field, Unshaded, Height field, Unshaded}
    \autofitgraphics[]{circular-without-radial-erosion-render.png, circular-with-radial-erosion-render.png}
    \autofitcaptions[]{Without radial drainage patterns, With radial drainage patterns}
    \caption[Radial drainage inserted in our dataset generation]{Multiple procedural strokes are added between the island centre and the surrounding sea, alternately landward and seaward. This transforms a simple circular island (left) into a more realistic volcanic island (right) by cheaply simulating star-shaped radial drainage patterns.}
    \label{fig:coral-island-radial-pattern-example}
\end{figure}

Once all inputs are set, we generate an example for multiple levels of subsidence $\subsidRate \in [0, 1]$ producing height fields that incorporate coral reef modelling along with its associated label map.

\begin{figure}[tb]
    \autofitgraphics[]{2_features.png, 2_heightmaps.png, 2_results.png}
    \autofitgraphics[]{1_features.png, 1_heightmaps.png, 1_results.png}
    \autofitcaptions{Input sketches, Output height field, Result}
    \caption[Multiple islands generated with two subsidence values]{An identical label map yields similar height fields over multiple inferences from the model, even after modifying the subsidence factor (visible in the luminosity of the input image).}
    \label{fig:coral-island-results-subsidence}
\end{figure}

% The pix2pix model was originally pretrained using RGB images. In this training phase, the images were labelled using the HSV (Hue, Saturation, Value) colour space, where the Hue component specifically carried the label information. Both the Saturation and Value components were kept neutral, meaning they did not convey any significant label-related data. The target images, the ones the model aimed to reproduce, were formatted in RGB.

For the purpose of training the pix2pix model with HSV images as input, we use the Hue component to encode the labels from the label map. We introduced a new dimension to the model's learning capabilities by incorporating the subsidence rate, denoted as $\subsidRate$, into the Value component.

Moreover, we purposefully left the Saturation component unchanged at this stage, reserving space for potentially including another parameter in the future, which would allow us to expand the model's utility without altering the foundational HSV encoding scheme established during its initial training. By adhering to this encoding format, we ensure continuity in data representation, which maximises the efficiency of the pretrained model. This strategic update enhances the model's adaptability and broadens its applicability to tackle new, complex challenges more effectively.

\begin{figure}[tb]
    \autofitgraphics[]{6_features.png, 6_heightmaps.png, 6_results.png}
    \autofitcaptions{Input sketch, Output height field, Result}
    \caption[A sample of the dataset with multiple transformations]{A sample of the dataset with the input and output of the island copied and scaled at multiple positions of the image, creating three different islands in a single image.}
    \label{fig:coral-island-cGAN-examples}
\end{figure}

This configuration enables rapid generation of large quantity of samples, with multiple parameters, of a single island centred in the image. We enhance each sample with several data augmentation techniques in addition to the usual affine transformations (rotation, scaling, and flipping).

\textit{Translation:} Since the original algorithm always centres the island, we translate the islands within the image to remove this constraint (\cref{fig:coral-island-cGAN-examples}). This ensures that the cGAN can generate islands in any position within the frame. By wrapping the island on the borders, the model learns that data can appear on the edges of the image.

\textit{Copy-paste:} We combine multiple islands into a single sample with only a maximum intersection over union (IoU) of 10\%, allowing the abysses to overlap but without obtaining other region types to merge. Height fields are blend using the smooth maximum function from \cref{eq:coral-island-smoothmax-function}, and label blending uses the usual $\max$ operator. Regions not covered by any island are assigned the abyss ID, which correspond to the minimal height.

All augmentation techniques are simultaneously applied to both the height field and the label map ensuring consistency between input (the label map) and output (the height field). 
% We summarise the complete dataset generation in \cref{alg:dataset-generation}.


% \begin{algorithm}
%     \caption{Procedural dataset generation.}
%     \label{alg:dataset-generation}
%     \DontPrintSemicolon
%     \KwIn{Base outline radii $\radius^*$, profile $\heightProfile^*$, resistance $\resistance^*$, noise $\noise$, stroke generator, subsidence grid $\{\subsidRate\}$}
%     \KwOut{Pairs $\{(\text{label map},\ \text{height field})\}$}

%     \BlankLine
%     \tcp{(1) Randomise sketches}
%     \ForEach{outline}{
%     $\radius(\angl) \gets \radius^* + \noise(\angl)$
%     }
%     $\heightProfile(\distRegions_\p) \gets \heightProfile^*(\distRegions_\p)\cdot \noise(\distRegions_\p)$\;
%     $\resistance(\distRegions_\p) \gets \resistance^*(\distRegions_\p)\cdot \noise(\distRegions_\p)$\;

%     \tcp{(2) Random strokes}
%     Generate $n$ strokes $\{\curve\}$ by uniformly sampling $m$ control points each; sample $\std$ and $S_\curve$ per stroke\;

%     \tcp{(3) Synthesis over subsidence levels}
%     \ForEach{$\subsidRate \in [0,1]$}{
%     Build $h$ and $I$ via \cref{alg:sketch-to-height-wind} and apply \cref{alg:subsidence-coral-blend} to get $\hat{h}$\;
%     label map $\gets$ HSV image with $I$ in the Hue channel and $\subsidRate$ in the Value channel\;
%     Store pair $(\text{label map},\, \hat{h})$\;
%     }

%     \tcp{(4) Augmentation}
%     Apply translations with wrapping, rotations, scalings, flips to both label map and height field consistently\;
%     Copy-paste multiple islands with IoU$\le 0.1$; blend heights with $\smoothmax$ and labels with $\max$; and assign abyss ID for uncovered regions\;

%     \Return dataset
% \end{algorithm}